%! Algebra 2 — Unit 1, Lesson 3 — Group Activity (Key)
\documentclass[10pt]{article}
\usepackage{algebra2-article}
\usepackage{algebra2-boxes}
\usepackage{algebra2-key}

\newcommand{\TallMath}[1]{$\displaystyle #1\rule[-1.4em]{0pt}{3.2em}$}

%=============================================================================
\begin{document}

\pageheader{Unit 1, Lesson 3}{Group Activity --- Key}
\namepartnerperiod

\vspace{0.15in}

% ─────────────────────────────────────────────
%  TIER R — Remediate
% ─────────────────────────────────────────────
\begin{tcolorbox}[colback=white,colframe=black!40,boxrule=0.7pt,arc=2mm,
    title={\bfseries Tier R — Remediate},fonttitle=\small,breakable]

\begin{enumerate}

    \item Function identification.

    \begin{enumerate}[label=(\alph*)]
        \item $\{(1,3),(2,5),(3,3),(4,7)\}$ \quad Answer: \ans{Yes} --- every input has exactly one output.
        \item $\{(2,4),(2,6),(3,8)\}$ \quad Answer: \ans{No} --- input $2$ maps to both $4$ and $6$.
    \end{enumerate}

    \begin{teachernote}
        Common error: students see the repeated output $3$ in (a) and call it not a function.
        Reinforce: outputs can repeat; inputs cannot.
    \end{teachernote}

    \vspace{0.1in}

    \item $f(x) = -2x + 8$
    \begin{enumerate}[label=(\alph*)]
        \item Slope: \ans{$-2$} --- the line falls 2 units for every 1 unit to the right.
        \item $y$-intercept: \ans{$(0, 8)$} \quad $x$-intercept (zero): \ans{$(4, 0)$} \; (set $f(x)=0$: $-2x+8=0 \Rightarrow x=4$)
    \end{enumerate}

    \vspace{0.1in}

    \item Parabola $f(x) = (x-1)^2 - 4$

    Vertex: \ans{$(1, -4)$} \quad Axis: \ans{$x = 1$} \quad $y$-int: \ans{$(0, -3)$} \quad Zeros: \ans{$x = -1$ and $x = 3$} \quad Opens: \ans{up}

    \begin{teachernote}
        Zeros: $(x-1)^2 = 4 \Rightarrow x-1 = \pm 2 \Rightarrow x = 3$ or $x = -1$.
    \end{teachernote}

    \vspace{0.1in}

    \item \ans{The trend appears linear because the data rises at an approximately constant rate (a straight line through the points).}

\end{enumerate}
\end{tcolorbox}

\vspace{0.15in}

% ─────────────────────────────────────────────
%  TIER A — Approaching Proficiency
% ─────────────────────────────────────────────
\begin{tcolorbox}[colback=white,colframe=black!40,boxrule=0.7pt,arc=2mm,
    title={\bfseries Tier A — Approaching Proficiency},fonttitle=\small,breakable]

\begin{enumerate}

    \item Line through $(-2, 1)$ and $(4, 13)$.

    \ansline{$m = \dfrac{13-1}{4-(-2)} = \dfrac{12}{6} = 2$}
    \ansline{$y - 1 = 2(x + 2) \Rightarrow y = 2x + 5$}
    Answer: \ans{$y = 2x + 5$}

    \vspace{0.1in}

    \item $f(x) = -x^2 + 6x - 5$
    \begin{enumerate}[label=(\alph*)]
        \item Axis of symmetry: \ans{$x = -\dfrac{6}{2(-1)} = 3$} \quad Vertex: \ans{$f(3) = -9+18-5 = 4$; vertex $(3, 4)$}
        \item \ansline{$-x^2 + 6x - 5 = 0 \Rightarrow x^2 - 6x + 5 = 0 \Rightarrow (x-1)(x-5) = 0$}
              Answer: \ans{$x = 1$ or $x = 5$}
        \item Domain: \ans{all real numbers} \quad Range: \ans{$(-\infty, 4]$} \; (opens down, max at $y=4$)
    \end{enumerate}

    \begin{teachernote}
        Students often forget to negate when converting $-x^2 + 6x - 5 = 0$ to standard form.
        Also watch for range written as $[4, \infty)$ --- the parabola opens \emph{down}, so values go to $-\infty$.
    \end{teachernote}

    \vspace{0.1in}

    \item $g(x) = 4 \cdot 2^x$

    \ansline{$g(0) = 4 \cdot 2^0 = 4 \cdot 1 = 4$}
    \ansline{$g(3) = 4 \cdot 2^3 = 4 \cdot 8 = 32$}
    Domain: \ans{all real numbers} \quad Range: \ans{$(0, \infty)$}

    \vspace{0.1in}

    \item $y = 3.5x + 12$ (subscribers context):

    \ans{Slope $3.5$: the number of subscribers increases by about $3{,}500$ per year.}\\
    \ans{$y$-intercept $12$: in 2010 (when $x=0$) there were approximately $12{,}000$ subscribers.}

\end{enumerate}
\end{tcolorbox}

% ─────────────────────────────────────────────
%  TIER E — Extension
% ─────────────────────────────────────────────
\begin{tcolorbox}[colback=white,colframe=black!40,boxrule=0.7pt,arc=2mm,
    title={\bfseries Tier E — Extension},fonttitle=\small,breakable]

\begin{enumerate}

    \item Perpendicular to $y = \dfrac{3}{4}x + 2$ through $(3, -1)$.

    \ansline{Perpendicular slope: $m = -\dfrac{4}{3}$}
    \ansline{$y - (-1) = -\dfrac{4}{3}(x - 3) \Rightarrow y + 1 = -\dfrac{4}{3}x + 4$}
    Answer: \ans{$y = -\dfrac{4}{3}x + 3$}

    \vspace{0.1in}

    \item Heart rate regression.
    \begin{enumerate}[label=(\alph*)]
        \item Desmos regression $y_1 \sim mx_1 + b$: \ans{approximately $y = -3.6x + 85.6$}
        \item \ans{Slope $-3.6$: each additional hour of exercise per week predicts a decrease of about $3.6$ bpm in resting heart rate.}
        \item \ans{At $x = 8$: $y \approx -3.6(8)+85.6 = 56.8$ bpm. This is within a reasonable biological range but is an extrapolation beyond the data, so reliability decreases.}
    \end{enumerate}

    \begin{teachernote}
        Regression coefficients may vary slightly depending on rounding in Desmos. Accept answers within $\pm 0.2$ of the slope.
        The key teaching point on (c) is evaluating validity of predictions beyond the data range.
    \end{teachernote}

    \vspace{0.1in}

    \item Comparison of $f(x)=2x+10$, $g(x)=x^2+2$, $h(x)=1.5^x$.

    \ansline{Graph all three in Desmos. $h(x) > f(x)$ for approximately $x > 9.4$ (intersection near $(9.4, 28.8)$).}
    \ans{For large $x$, the exponential $h(x)$ eventually dominates the linear $f(x)$, but it takes until $x \approx 9.4$.}

\end{enumerate}
\end{tcolorbox}

\end{document}
